\chapter{Hyperbaric Oxygen Therapy}

\section*{Overview}
Hyperbaric oxygen therapy (HBOT) is a treatment in which the patient breathes 100\% oxygen while inside a pressurized chamber at greater than atmospheric pressure. It is used to treat conditions involving tissue hypoxia, gas bubbles, or infection.

\section*{Alternative Names}
Hyperbaric oxygen therapy (HBOT); hyperbaric medicine; high-pressure oxygen therapy; HBO

\section*{Principle}
Breathing 100\% oxygen under increased ambient pressure markedly raises the partial pressure of oxygen dissolved in plasma and delivered to tissues, independent of hemoglobin. This corrects tissue hypoxia, promotes angiogenesis and wound healing, exerts antimicrobial effects, and shrinks gas bubbles in decompression illness by creating a pressure gradient for gas reabsorption.

\section*{History}
Hyperbaric oxygen use dates to the 19th century, and after the success of compressed-air chambers in diving medicine, its modern clinical applications were developed in the mid-20th century.

\section*{Why It Is Done}
Ordered for decompression sickness, arterial gas embolism, carbon monoxide poisoning (especially with neurologic signs or myocardial injury), and radiation injury to bone and soft tissue. It is also used for chronic nonhealing diabetic wounds, necrotizing soft-tissue infections, crush injuries, compromised skin grafts, and certain refractory osteomyelitis.

\section*{Is This a Primary or Secondary Test or Procedure}
A primary treatment for decompression sickness, arterial gas embolism, and severe carbon monoxide poisoning, and an adjunctive therapy for most other indications.

\section*{How It Is Performed}
The patient enters a monoplace chamber or a multiplace chamber where a technician controls pressure. The chamber is pressurized to approximately 2 to 3 atmospheres absolute while the patient breathes 100\% oxygen, typically for sessions of 60 to 120 minutes. Multiple daily sessions are usually required over several weeks.

\section*{Preparation}
The patient should have no recent alcohol use and should avoid smoking during treatment. A pretreatment assessment includes a chest radiograph or pulmonary evaluation and oximetry, and the patient wears only cotton clothing to reduce the fire risk of oxygen-rich environments.

\section*{Contraindications and Precautions}
The only absolute contraindication is untreated pneumothorax, because trapped gas expands on ascent. Relative contraindications include COPD with bullae, active infections, severe cardiac disease, pregnancy, claustrophobia, and current use of drugs such as bleomycin or doxorubicin that potentiate oxygen toxicity.

\section*{Risks}
Common adverse effects are reversible middle-ear or sinus barotrauma during compression and temporary myopia. Serious but rare complications include oxygen seizures, pulmonary oxygen toxicity, and pneumothorax, and a brief risk of fire exists within the oxygen-rich chamber.

\section*{Aftercare and Recovery}
The patient is observed briefly after each session for ear discomfort or dizziness and instructed on pressure-equalization techniques. The course continues with serial assessment of the wound or underlying condition to determine the number of sessions needed.

\section*{Outcome and Follow-up}
Response is assessed by wound healing, resolution of neurologic signs, or improvement in imaging and clinical status. The number of sessions is adjusted based on these endpoints.

\section*{Expected Results}
Healing of the target wound, resolution of gas bubble or carbon monoxide effects, and recovery of neurologic or tissue function within the expected treatment course.

\section*{Follow-up}
Wound care continues with regular measurement and photography, and serial examinations document neurologic recovery. When the course is complete, the patient is discharged with instructions to avoid flying or high-altitude exposure until cleared, and long-term follow-up addresses any residual deficits.

\section*{When to Seek Medical Care}
Follow the procedure team's specific instructions. Seek urgent medical assessment for severe or worsening pain, uncontrolled bleeding, fever or signs of infection, trouble breathing, chest pain, fainting, new weakness or confusion, inability to urinate, or any rapidly worsening symptom. The appropriate warning signs depend on the procedure and the patient's medical condition.

\section*{References}
\begin{enumerate}
  \item MedlinePlus. Hyperbaric oxygen therapy. U.S. National Library of Medicine; 2025.
  \item Current Medical Diagnosis and Treatment. McGraw-Hill; 2025.
\end{enumerate}

\clearpage
